% !TeX root = ../main.tex
\chapter{Objectives}\label{ch:objectives}

\begin{guidance}
State the general objective of the thesis from a high-level perspective, and break it
down into specific, measurable objectives. A reader should get a clear idea of what has
been done just by reading this chapter. If your work is hypothesis-driven, state the
hypothesis or research questions first. The specific objectives will be revisited in
\autoref{ch:conclusions} to assess whether they have been achieved.
\end{guidance}

\section{Research questions}\label{sec:research-questions}

\begin{enumerate}
  \renewcommand{\labelenumi}{\textbf{RQ\arabic{enumi}.}}
  \renewcommand{\theenumi}{RQ\arabic{enumi}}
  \item\label{rq:1} [First research question]
  \item\label{rq:2} [Second research question]
\end{enumerate}

\section{General objective}\label{sec:general-objective}

[One sentence describing the overall aim of the thesis.]

\section{Specific objectives}\label{sec:specific-objectives}

\begin{enumerate}
  \renewcommand{\labelenumi}{\textbf{O\arabic{enumi}.}}
  \renewcommand{\theenumi}{O\arabic{enumi}}
  \item\label{obj:1} [Review the state of the art on \dots]
  \item\label{obj:2} [Collect and prepare the data \dots]
  \item\label{obj:3} [Design and implement \dots]
  \item\label{obj:4} [Evaluate \dots and compare with \dots]
\end{enumerate}

\section{Scope and limitations}\label{sec:scope}

[Define what is within the scope of the thesis and what is deliberately left out.]
